\chapter{HMC for Occasionally-Binding Constraints with Discontinuous Gradients}
\label{ch:obc-discontinuous-gradient-hmc}

\section{Motivation and Scope}

Occasionally-binding constraint (OBC) models create a fundamental challenge for
Hamiltonian Monte Carlo when the constraint enforcement mechanism induces
discontinuities in the marginal likelihood gradient. This chapter addresses the
specific case where:
\begin{itemize}
  \item An OBC solver selects regime paths based on equilibrium consistency
    conditions;
  \item The selection can be discrete (integer solver) or smoothed (softplus,
    logistic);
  \item A nonlinear filter (UKF, particle filter) computes the marginal
    likelihood $p(y_{1:T} \mid \theta)$;
  \item The composed map $\theta \mapsto \log p(y \mid \theta)$ has regions of
    steep or discontinuous gradient variation;
  \item Standard leapfrog HMC suffers from poor acceptance or stuck chains.
\end{itemize}

The zero lower bound (ZLB) on nominal interest rates is the canonical
application. When the shadow rate falls below zero, the observed rate is
censored, and different policy regimes (quantitative easing on/off) can create
discrete switches in the model's equilibrium dynamics. The filtering
distribution can collapse near the bound, amplifying parameter-space gradient
discontinuities.

This chapter does not present a complete solution. It maps the problem
geometry, surveys relevant HMC extensions, and proposes a research program for
making gradient-based MCMC tractable in this setting. The focus is on
\emph{practical inference for finite-horizon DSGE models}, not on existence
theorems or asymptotic optimality.

\section{The Composed Likelihood Map}

\subsection{Three Layers of Computation}

The parameter-to-likelihood map for an OBC state-space model has three
distinct layers:
\begin{equation}
  (\theta, x_{t-1}, \varepsilon_t)
  \xrightarrow{\mathcal{S}} (r_t, F_{r_t}, Q_{r_t}, H_{r_t})
  \xrightarrow{\mathcal{F}} \log p(y_{1:T} \mid \theta)
  \xrightarrow{\mathcal{K}_{\text{HMC}}} \theta'.
  \label{eq:obc-three-layers}
\end{equation}

\paragraph{Layer 1: OBC solver $\mathcal{S}$.}
Maps parameters and a candidate state/shock trajectory to:
\begin{itemize}
  \item A regime path $r_1, \ldots, r_T$ over a finite horizon;
  \item Regime-dependent state-space matrices $(F_{r_t}, Q_{r_t}, H_{r_t})$ at
    each time $t$.
\end{itemize}
The solver can be:
\begin{description}
  \item[Integer/hard] Selects the first regime path satisfying all equilibrium
    consistency inequalities (OccBin, Boehl spell-duration, PyDSGE affine
    step). Creates discrete regime switches.
  \item[Smoothed] Uses softplus, logistic, or weighted averaging to
    interpolate between regimes. Creates continuous but steep gradients.
\end{description}

\paragraph{Layer 2: Filter $\mathcal{F}$.}
Computes the marginal log-likelihood:
\begin{equation}
  \log p(y_{1:T} \mid \theta)
  = \sum_{t=1}^T \log p(y_t \mid y_{1:t-1}, \theta),
\end{equation}
using the regime-dependent matrices from the solver. Common choices:
\begin{itemize}
  \item Kalman filter (piecewise linear model);
  \item Extended Kalman filter (first-order Taylor);
  \item Unscented/sigma-point filter (deterministic sampling);
  \item Particle filter (stochastic Monte Carlo).
\end{itemize}

\paragraph{Layer 3: HMC kernel $\mathcal{K}_{\text{HMC}}$.}
Proposes $\theta' \sim \mathcal{K}(\theta' \mid \theta)$ using Hamiltonian
dynamics driven by $\nabla_\theta \log p(y \mid \theta)$ and accepts via
Metropolis correction. Requires:
\begin{itemize}
  \item Evaluable gradient almost everywhere;
  \item Reversible, volume-preserving proposal map;
  \item Bounded integration error or exact event handling.
\end{itemize}

\subsection{Where Nonsmoothness Enters}

The gradient $\nabla_\theta \log p(y \mid \theta)$ is a composition:
\begin{equation}
  \nabla_\theta \log p(y \mid \theta)
  = \nabla_\theta \left[ \mathcal{F} \circ \mathcal{S}(\theta) \right].
  \label{eq:gradient-composition}
\end{equation}

Discontinuities can arise from:
\begin{enumerate}
  \item \textbf{Solver branch switches.} An integer solver changes the regime
    path when a margin inequality flips. The state-space matrices jump
    discretely, even if $\theta$ changes smoothly.
  \item \textbf{Filter collapse.} When observations pin a state variable to a
    constraint boundary (e.g., observed rate at ZLB for many periods), the
    filtering variance collapses. Small $\theta$ perturbations can shift the
    entire sigma-point cloud across the boundary, causing sharp likelihood
    changes.
  \item \textbf{Implicit boundaries.} The switching surface
    $\{\theta : g(\theta) = 0\}$ is not given analytically. It is defined
    implicitly by the solver's equilibrium conditions, which depend on future
    state paths that themselves depend on earlier regime choices.
\end{enumerate}

The result is a likelihood surface with:
\begin{itemize}
  \item Piecewise-smooth regions (branch interiors);
  \item Boundaries where the gradient is undefined or jumps;
  \item Possible value discontinuities if branch likelihoods do not match.
\end{itemize}

\section{Geometry Classification}

\subsection{Kink versus Value Jump}

Let $\mathcal{R}_j$ denote a parameter region with a fixed regime path, and
let $\ell_j(\theta) = \log p(y \mid \theta, \text{regime } j)$ be the
branch-interior log-likelihood. At a switching boundary $\theta_b$, there are
three cases:

\paragraph{Case 1: Smooth boundary (no issue).}
Values and gradients match:
\begin{equation}
  \lim_{\theta \to \theta_b, \theta \in \mathcal{R}_j} \ell_j(\theta)
  = \lim_{\theta \to \theta_b, \theta \in \mathcal{R}_k} \ell_k(\theta),
  \quad
  \lim_{\theta \to \theta_b, \theta \in \mathcal{R}_j} \nabla \ell_j(\theta)
  = \lim_{\theta \to \theta_b, \theta \in \mathcal{R}_k} \nabla \ell_k(\theta).
\end{equation}
Standard HMC applies without modification.

\paragraph{Case 2: Continuous kink.}
Values match but gradients do not:
\begin{equation}
  \lim_{\theta \to \theta_b, \theta \in \mathcal{R}_j} \ell_j(\theta)
  = \lim_{\theta \to \theta_b, \theta \in \mathcal{R}_k} \ell_k(\theta),
  \quad
  \lim_{\theta \to \theta_b, \theta \in \mathcal{R}_j} \nabla \ell_j(\theta)
  \ne \lim_{\theta \to \theta_b, \theta \in \mathcal{R}_k} \nabla \ell_k(\theta).
\end{equation}
The target density is continuous, but the gradient has a jump discontinuity.
Standard leapfrog remains \emph{valid} (no reflection needed), but
\emph{performs poorly} due to integration error at crossings.

\paragraph{Case 3: Value jump.}
Values do not match:
\begin{equation}
  \lim_{\theta \to \theta_b, \theta \in \mathcal{R}_j} \ell_j(\theta)
  \ne \lim_{\theta \to \theta_b, \theta \in \mathcal{R}_k} \ell_k(\theta).
\end{equation}
The potential $U(\theta) = -\log p(\theta \mid y)$ has a finite jump
$\Delta U$. Standard leapfrog is \emph{invalid}: the trajectory cannot
conserve energy across the jump without an explicit event map (reflection or
refraction).

\subsection{Why Standard Leapfrog Fails}

For a smooth target with potential $U(\theta)$ and quadratic kinetic energy
$K(p) = \frac{1}{2} p^\top M^{-1} p$, the leapfrog integrator is:
\begin{equation}
  p_{1/2} = p_0 - \frac{\epsilon}{2} \nabla U(\theta_0),
  \quad
  \theta_1 = \theta_0 + \epsilon M^{-1} p_{1/2},
  \quad
  p_1 = p_{1/2} - \frac{\epsilon}{2} \nabla U(\theta_1).
  \label{eq:standard-leapfrog}
\end{equation}
The Hamiltonian error is $O(\epsilon^2)$ under smoothness. Energy conservation
and volume preservation ensure the Metropolis correction targets the correct
posterior.

\paragraph{Failure at a continuous kink.}
When $\nabla U$ has a jump, a trajectory that crosses the boundary evaluates
the gradient on one side at $\theta_0$ and on the other side at $\theta_1$.
The integrator does not ``see'' the jump itself—it only samples the gradient
at discrete points. The energy error is $O(1)$ in the jump size, not $O(\epsilon^2)$.
As $\epsilon \to 0$, the crossing is localized more accurately, but the
missing gradient contribution remains finite.

Result: large Hamiltonian error $\to$ low acceptance $\to$ chain stuck near
boundaries or in one regime.

\paragraph{Failure at a value jump.}
A leapfrog step that crosses a potential discontinuity pays no energy cost for
the jump. If $\Delta U > 0$ and the trajectory moves from the low-potential to
high-potential region, the proposal systematically underestimates the true
Hamiltonian. Acceptance rates collapse, and the chain cannot cross the
boundary reversibly.

\subsection{OBC-Specific Complications}

\paragraph{Implicit boundaries.}
In a finite-horizon DSGE model, the regime at time $t$ depends on:
\begin{itemize}
  \item Parameters $\theta$;
  \item Lagged states (which depend on regimes $r_{1:t-1}$);
  \item Future equilibrium conditions (which depend on regimes $r_{t+1:T}$).
\end{itemize}
The switching surface $g(\theta) = 0$ is not a simple hyperplane. It is an
implicit, high-dimensional manifold defined by the solver's consistency check.

\paragraph{Filter amplification.}
A solver coefficient jump can be small, but the filtered likelihood change can
be large. The innovation covariance $S_t$, Kalman gain $K_t$, and log
determinant $\log |S_t|$ all depend on the regime-dependent matrices. A
boundary crossing that changes $F_t$ or $H_t$ can produce a large
$\Delta \ell$ even when the raw margin change is small.

\paragraph{Shadow-rate identification collapse.}
When the observed rate is at the ZLB for multiple periods, the filtering
distribution over the shadow rate collapses. All deeply-negative shadow values
produce the same bounded observed rate (up to small observation noise). The
likelihood provides no information about how far below zero the shadow rate is.

This creates a \emph{plateau} in parameter space: $\theta$ values that differ
only in their implied shadow-rate depth have nearly identical likelihoods. HMC
trajectories can wander on this plateau with little force directing them, even
though the gradient is well-defined (zero) almost everywhere.

\section{Existing HMC Methods for Discontinuous Targets}

\subsection{Reflection and Refraction HMC}

\subsubsection{Continuous Hamiltonian flow with event detection}

For a potential with a jump $\Delta U$ across a codimension-one surface
$g(\theta) = 0$, Afshar \& Domke (2015) and Nishimura, Dunson, \& Lu (2020)
derive an event map that preserves the Hamiltonian exactly. Let $\nu = \nabla
g / \|\nabla g\|$ be the outward normal. At the event time $t_e$:
\begin{enumerate}
  \item Decompose momentum: $p = p_n \nu + p_{\text{tan}}$, where
    $p_n = \nu^\top p$.
  \item Compute normal kinetic energy:
    $K_n = p_n^2 / (2 m_n)$, where $m_n^{-1} = \nu^\top M^{-1} \nu$.
  \item If $K_n < \Delta U$: reflect, $p^+ = p - 2 p_n \nu$.
  \item If $K_n \geq \Delta U$: refract,
    \begin{equation}
      p_n^+ = \operatorname{sign}(p_n^-) \sqrt{(p_n^-)^2 - 2 m_n \Delta U},
      \quad
      p^+ = p^- + (p_n^+ - p_n^-) \nu.
      \label{eq:refraction-rule}
    \end{equation}
\end{enumerate}

\paragraph{Why this preserves the target.}
The event map is a deterministic limit of smooth mollified flows. It conserves
energy: $U(\theta^+) + K(p^+) = U(\theta^-) + K(p^-)$. It preserves
phase-space volume. The reverse map (flip momentum, apply reverse potential
jump) recovers the incoming state. Under stated regularity conditions, the
composed flow (smooth segments + event maps) is a valid HMC proposal.

\paragraph{Required information.}
\begin{enumerate}
  \item Boundary function $g(\theta)$ and its gradient $\nabla g(\theta)$.
  \item Event detection: locate the first crossing time along the trajectory.
  \item One-sided potential values: $U_j(\theta_e)$ and $U_k(\theta_e)$ to
    compute $\Delta U$.
  \item Momentum metric: mass matrix $M$ to compute $m_n$.
\end{enumerate}

\paragraph{OBC implementation burden.}
For an implicit OBC boundary:
\begin{itemize}
  \item $g(\theta)$ is not given analytically. It is defined by the solver's
    equilibrium check, which requires full solution and filtering.
  \item Detecting a crossing requires bisection or root-finding, re-solving
    the model at intermediate $\theta$ values.
  \item The event normal $\nu$ is a total derivative through the solver $\to$
    filter composition, not a simple hyperplane normal.
  \item Multiple margins may be active simultaneously; a robust tie rule is
    needed.
\end{itemize}

This is computationally expensive but mathematically correct for discrete
regime switches with value jumps.

\subsubsection{Discontinuous HMC with Laplace momentum}

Nishimura et al. (2020) introduce a coordinate-wise integrator for targets
with discontinuous coordinates. For a coordinate $\theta_i$ whose potential
has jumps:
\begin{enumerate}
  \item Use Laplace momentum: $K(p_i) = |p_i| / m_i$ (constant-speed motion).
  \item Move the coordinate at constant velocity until hitting a boundary or
    exhausting the time budget.
  \item At each boundary, check if available kinetic energy can pay the
    potential change. If yes, cross; if no, reflect.
  \item Use a random coordinate order (symmetric under reversal) to ensure
    reversibility in distribution.
  \item Use a randomized step size to avoid lattice artifacts.
\end{enumerate}

\paragraph{Advantages.}
\begin{itemize}
  \item No need to locate exact event times in continuous trajectories.
  \item Can handle multiple discontinuities in one coordinate move.
  \item Computationally simpler than continuous-trajectory event detection.
\end{itemize}

\paragraph{Disadvantages.}
\begin{itemize}
  \item Random coordinate order and step size are \emph{required}, not
    optional. A fixed step can produce a chain that moves on a grid, failing
    irreducibility even when visually active.
  \item Less efficient for strongly correlated posteriors (coordinate-wise
    moves vs coherent trajectories).
  \item Still requires evaluating the potential on both sides of each
    discontinuity.
\end{itemize}

\paragraph{Mixed construction for OBC.}
Use DHMC-style Laplace coordinates for parameters identified as
regime-sensitive (those that move margins), and ordinary Gaussian HMC for
smooth coordinates. This requires:
\begin{itemize}
  \item A parameter-sensitivity analysis to identify which $\theta_i$ affect
    regime switches;
  \item Stable hard-potential evaluation after each coordinate move;
  \item Adaptation rules compatible with the mixed kernel.
\end{itemize}

\subsection{Generalized Randomized HMC for Piecewise-Smooth Targets}

Tran \& Kleppe (2025) address the \emph{continuous kink} case most directly.
Their Generalized Randomized HMC (GRHMC):
\begin{enumerate}
  \item Integrates the Hamiltonian flow inside each smooth region using
    Runge--Kutta or leapfrog.
  \item When a numerical step would cross a boundary, truncate at the first
    event.
  \item Update the gradient to the new region's value.
  \item Continue the remaining part of the step under the new gradient.
  \item The position is \emph{not} jumped—only the force changes.
\end{enumerate}

For a genuinely piecewise-density target (value jump), they combine this with
boundary reflection/refraction.

\paragraph{Applicability to OBC.}
If the OBC filtered likelihood is continuous with discontinuous gradient (Case
2), GRHMC is a direct candidate. The implementation must:
\begin{itemize}
  \item Detect when a leapfrog step crosses an implicit solver margin
    (bisection + re-solve);
  \item Recompute the filtered likelihood and gradient at the event;
  \item Handle multiple or grazing events.
\end{itemize}

If the likelihood has value jumps (Case 3), the piecewise-density
reflection/refraction extension is needed.

\paragraph{What this does not solve.}
GRHMC assumes the boundary geometry can be queried. For an implicit
finite-horizon OBC solver, locating $g(\theta) = 0$ and computing $\nabla
g(\theta)$ requires full model solution at intermediate points. The method is
mathematically valid but computationally intensive.

\subsection{Kernel and Neural-Proposal HMC}

Strathmann et al. (2015) propose Kernel Hamiltonian Monte Carlo (KMC): fit a
surrogate potential $\widetilde{U}(\theta)$ using score matching, use its
gradient to drive proposals, but evaluate the \emph{true} target $U(\theta)$
in the Metropolis correction:
\begin{equation}
  \alpha = \min\left\{1, \exp\left[ -U(\theta^*) - K(p^*) + U(\theta) + K(p) \right] \right\}.
  \label{eq:kmc-correction}
\end{equation}

The surrogate changes proposal geometry but not the stationary target.

\paragraph{Neural proposal for OBC.}
Train a neural network $s_\phi(\theta)$ to approximate $\nabla_\theta \log
p(y \mid \theta)$ using:
\begin{itemize}
  \item Branch-interior gradients (automatic differentiation through the hard
    solver + filter);
  \item Points near each active margin;
  \item Explicit cross-boundary pairs to capture the gradient jump;
  \item Lipschitz or Sobolev regularization for smooth interpolation.
\end{itemize}

Use $s_\phi$ to drive the leapfrog proposal, but evaluate the hard likelihood
for acceptance. A poor surrogate leads to low acceptance, not wrong posterior.

\paragraph{Adaptation and exactness.}
If the surrogate is refit from chain history, the kernel is adaptive.
Adaptation must either:
\begin{itemize}
  \item Diminish over time (Roberts \& Rosenthal conditions);
  \item Freeze after a finite warm-up phase before collecting posterior
    samples.
\end{itemize}

A held-out gradient loss is a training diagnostic. It does not prove posterior
correctness—only the hard-target acceptance ratio does that.

\section{Smoothing Strategies and Their Targets}

\subsection{Four Distinct Smoothing Operations}

``Smoothing'' can refer to four mathematically different operations:

\paragraph{1. Economic policy rule.}
Replace a hard constraint $r = \max\{r^{\text{shadow}}, \bar{r}\}$ with a
smooth approximation:
\begin{equation}
  r = \bar{r} + \operatorname{softplus}(r^{\text{shadow}} - \bar{r}, \alpha)
  = \bar{r} + \frac{1}{\alpha} \log(1 + \exp(\alpha (r^{\text{shadow}} - \bar{r}))).
  \label{eq:softplus-rule}
\end{equation}

This changes the equilibrium model itself. As $\alpha \to \infty$, the
softplus approaches the hard max, but at finite $\alpha$ it defines a
different economic model. This is a legitimate modeling choice if the
smooth constraint better represents actual policy, but it is not a
computational trick to preserve the hard-max posterior.

\paragraph{2. Regime selector smoothing.}
Keep the economic model but replace a discrete regime choice with a weighted
average:
\begin{equation}
  F_{\text{smooth}} = (1 - w) F_{\text{regime 0}} + w F_{\text{regime 1}},
  \quad w = \operatorname{logistic}(g(\theta) / \tau).
  \label{eq:soft-selector}
\end{equation}

This creates a smooth map $\theta \mapsto F_{\text{smooth}}(\theta)$, but it
does not correspond to any single equilibrium. It mixes transition matrices
from two different models. Unless this mixture has a probabilistic
interpretation (latent regime model), it is neither the hard-selector
posterior nor a valid alternative model.

\paragraph{3. Filter smoothing.}
Smooth the filtering recursion or the final likelihood. For example, average
likelihoods from different sigma points that select different regimes, or
mollify the likelihood surface in parameter space:
\begin{equation}
  \ell_\delta(\theta) = \int \ell_{\text{hard}}(\eta) \, \varphi_\delta(\theta - \eta) \, \mathrm{d}\eta.
  \label{eq:mollified-likelihood}
\end{equation}

This targets a different posterior:
$\pi_\delta(\theta \mid y) \propto \exp\{\ell_\delta(\theta)\} \pi_0(\theta)$.
It can be studied as $\delta \to 0$, but at finite $\delta$ it is not the
hard posterior.

\paragraph{4. HMC proposal smoothing (preserves target).}
Keep the hard likelihood for acceptance but use a smooth surrogate for the
proposal gradient:
\begin{equation}
  \text{Proposal: } \theta^* = \Phi_{\widetilde{\nabla}}(\theta, p),
  \quad
  \text{Accept: } \alpha = \min\left\{1, \frac{p(\theta^* \mid y)}{p(\theta \mid y)} \frac{\varphi(p^*)}{(\varphi(p))} \right\},
  \label{eq:surrogate-proposal}
\end{equation}
where the acceptance uses the hard $p(\theta \mid y)$.

Only this fourth operation automatically preserves the hard posterior. The
first three define modified models or posteriors and must be labeled as such.

\subsection{Softplus as a Continuous-Kink Model}

Using softplus for the economic constraint (operation 1) creates:
\begin{itemize}
  \item A continuous observation function $h(\theta)$ with no value jumps;
  \item Steep but continuous gradients near the effective bound;
  \item No discrete regime switches (assuming no separate integer QE selector).
\end{itemize}

This is a \emph{continuous-kink geometry} (Case 2), not a value-jump geometry
(Case 3). Standard leapfrog remains valid but performs poorly when:
\begin{enumerate}
  \item The softplus temperature $\alpha$ is large (approaching hard max);
  \item The UKF filtering variance collapses near the bound;
  \item Small $\theta$ perturbations create large gradient changes.
\end{enumerate}

The gradient Lipschitz constant grows with $\alpha$ and with the degree of
filter collapse. As $\alpha \to \infty$, the softplus converges to the hard
max, and the gradient approaches a true discontinuity.

\subsection{Measurement Protocol for OBC Geometry}

Before choosing an HMC method, measure the actual boundary geometry:

\paragraph{Step 1: Value continuity check.}
For $\theta$ pairs straddling a detected regime switch:
\begin{equation}
  \Delta \ell = |\ell_j(\theta^-) - \ell_k(\theta^+)|, \quad \theta^+ - \theta^- = \epsilon \nu,
  \label{eq:value-jump-check}
\end{equation}
where $\nu$ is the boundary normal direction and $\epsilon \to 0$.

If $\Delta \ell / |\ell_j| < \text{tol}$: continuous kink (Case 2).

If $\Delta \ell / |\ell_j| = O(1)$: value jump (Case 3).

\paragraph{Step 2: Gradient jump magnitude.}
For continuous-kink boundaries:
\begin{equation}
  \Delta g = \|\nabla \ell_j(\theta^-) - \nabla \ell_k(\theta^+)\|_2.
  \label{eq:gradient-jump-check}
\end{equation}

Large $\Delta g$ relative to typical $\|\nabla \ell\|$ indicates poor
leapfrog geometry.

\paragraph{Step 3: Boundary crossing rate.}
Run short HMC trajectories and record:
\begin{itemize}
  \item Fraction of trajectories that cross a boundary;
  \item Number of crossings per trajectory;
  \item Acceptance rate conditional on crossing vs not crossing.
\end{itemize}

High crossing rate + low crossing acceptance $\to$ boundary is a bottleneck.

\paragraph{Step 4: Gradient Lipschitz estimation.}
Sample $\theta$ pairs with $\|\theta_1 - \theta_2\| = \delta$:
\begin{equation}
  L_{\text{est}} = \max_{\text{pairs}} \frac{\|\nabla \ell(\theta_1) - \nabla \ell(\theta_2)\|}{\|\theta_1 - \theta_2\|}.
  \label{eq:lipschitz-estimate}
\end{equation}

For softplus with temperature $\alpha$, $L$ grows with $\alpha$.

\paragraph{Step 5: Filter variance at bound.}
For periods where observed rate is at or near the ZLB, record the UKF
filtering variance of the shadow-rate state. Variance collapse (orders of
magnitude smaller than prior) indicates identification failure and gradient
plateau.

\section{Proposed Research Program}

\subsection{Phase 0: Documentation and Literature Review}

\begin{itemize}
  \item \textbf{Write this chapter} documenting the problem, geometry
    classification, and existing methods.
  \item Survey HMC extensions: reflection/refraction (Afshar \& Domke,
    Nishimura et al.), GRHMC (Tran \& Kleppe), KMC (Strathmann et al.).
  \item Document the BGS/ZLB-specific complications: implicit boundaries,
    finite-horizon solver, filter collapse, shadow-rate plateaus.
  \item Establish the decision tree: value jump $\to$ event-aware HMC or
    DHMC; continuous kink $\to$ GRHMC or surrogate proposal.
\end{itemize}

\subsection{Phase 1: Geometry Characterization}

\textbf{Target:} Two-country ZLB DNS model, softplus-smoothed bound, UKF
filtering.

\textbf{Experiments:}
\begin{enumerate}
  \item Vary softplus temperature $\alpha \in \{5, 10, 20, 50, 100\}$ and
    measure gradient Lipschitz constant.
  \item Run standard HMC at each $\alpha$ with fixed step size and record
    acceptance rate, ESS, boundary-crossing rate.
  \item Document filter variance collapse: plot UKF variance of shadow level
    vs time for periods at ZLB.
  \item Measure gradient jump magnitude $\Delta g$ near effective boundaries
    (where softplus is steep).
  \item Check value continuity: is $\ell(\theta)$ continuous or does softplus
    still produce small jumps through filter composition?
\end{enumerate}

\textbf{Deliverable:} A geometry report classifying the softplus-UKF target as
continuous-kink with quantified gradient variation and crossing-rate
diagnostics.

\subsection{Phase 2: Event-Aware Integration}

\textbf{Method:} Implement Tran--Kleppe-style truncated leapfrog.

\textbf{Algorithm:}
\begin{enumerate}
  \item Define a steep-gradient detector: flag when $\|\nabla \ell(\theta)\|$
    exceeds a threshold or when a monitored margin is near zero.
  \item Within a leapfrog step, if the detector flags a boundary:
    \begin{itemize}
      \item Bisect to locate the boundary crossing;
      \item Re-solve the OBC model and re-filter at the boundary $\theta_b$;
      \item Recompute $\nabla \ell(\theta_b)$ on the new side;
      \item Continue the step with the updated gradient.
    \end{itemize}
  \item If no boundary is flagged, use standard leapfrog.
\end{enumerate}

\textbf{Experiments:}
\begin{enumerate}
  \item Compare standard leapfrog vs event-aware leapfrog at $\alpha = 20$
    (moderate smoothing).
  \item Measure: acceptance rate, ESS per gradient evaluation, cost per step
    (including bisection overhead).
  \item Validate reversibility: record a trajectory, reverse momentum, replay,
    check $\|\theta_{\text{final}} - \theta_{\text{initial}}\| < \epsilon$.
\end{enumerate}

\textbf{Deliverable:} Cost-benefit analysis: does event detection improve
mixing enough to justify the re-solve overhead?

\subsection{Phase 3: Surrogate Gradient Proposal}

\textbf{Method:} Train a neural network $s_\phi(\theta)$ to approximate
$\nabla_\theta \log p(y \mid \theta)$.

\textbf{Training data:}
\begin{enumerate}
  \item Sample $\theta$ from the prior and from a short HMC warm-up.
  \item Compute hard $\nabla \ell(\theta)$ via automatic differentiation
    through the softplus solver + UKF.
  \item For points near boundaries (detected by margin proximity), oversample
    to ensure the network learns the steep region.
  \item Include explicit cross-boundary pairs $(\theta^-, \theta^+)$ to
    capture gradient variation.
\end{enumerate}

\textbf{Network architecture:}
\begin{itemize}
  \item Input: $\theta \in \mathbb{R}^d$ (optionally augment with margin
    telemetry).
  \item Output: $s_\phi(\theta) \in \mathbb{R}^d$ (approximate score).
  \item Loss: $\mathcal{L} = \mathbb{E}\left[ \|s_\phi(\theta) - \nabla \ell(\theta)\|^2 \right]$
    plus Lipschitz or Sobolev regularization.
\end{itemize}

\textbf{HMC with surrogate:}
\begin{enumerate}
  \item Use $s_\phi(\theta)$ to drive leapfrog proposals.
  \item Evaluate the \emph{hard} $\ell(\theta)$ and $\ell(\theta^*)$ for
    Metropolis acceptance.
  \item Freeze the network after warm-up before collecting posterior samples.
\end{enumerate}

\textbf{Experiments:}
\begin{enumerate}
  \item Compare: standard HMC (hard gradient), surrogate HMC (neural gradient
    + hard acceptance), fully-smoothed HMC (neural gradient + neural
    acceptance).
  \item Measure: acceptance rate, ESS per hard likelihood evaluation,
    posterior agreement (KL divergence, marginal CI coverage).
  \item Validate: a poor surrogate should produce low acceptance but correct
    posterior; a good surrogate should produce high acceptance and correct
    posterior.
\end{enumerate}

\textbf{Deliverable:} Evidence that neural-proposal HMC can preserve the hard
posterior while improving acceptance, or identification of failure modes
(boundary overfitting, tail errors).

\subsection{Phase 4: UKF Variance Inflation}

\textbf{Hypothesis:} Filter variance collapse amplifies gradient
discontinuities. Artificially inflating UKF variance near bounds may smooth
the marginal likelihood gradient.

\textbf{Method:}
\begin{enumerate}
  \item Detect when a state (e.g., shadow level) is near a bound and the
    filtering variance is below a threshold.
  \item Inflate the predicted covariance: $P_{t|t-1}^+ = P_{t|t-1} + \delta I$,
    where $\delta$ is calibrated to prevent collapse.
  \item Continue UKF update with the inflated covariance.
  \item The modified filter changes the likelihood, so this is \emph{not} a
    proposal-only trick—it defines a modified posterior.
\end{enumerate}

\textbf{Experiments:}
\begin{enumerate}
  \item Measure gradient Lipschitz with and without variance inflation.
  \item Run HMC on the inflated-filter likelihood and check: does it mix
    better? Does the posterior remain scientifically meaningful?
  \item Compare to: particle filter with maintained diversity (if/when
    gradients are available).
\end{enumerate}

\textbf{Deliverable:} Assessment of whether variance inflation is a practical
compromise (easier than particle filter, better than collapsed UKF) or
scientifically invalid (distorts posterior too much).

\subsection{Phase 5: Hard OBC with Integer Regime Selector}

\textbf{Target:} Full ZLB model with QE on/off regimes, integer solver.

\textbf{Geometry:} Likely value jumps (Case 3), not just continuous kinks.

\textbf{Method:} Reflection/refraction HMC (Nishimura et al. or Tran--Kleppe
piecewise-density extension).

\textbf{Required infrastructure:}
\begin{enumerate}
  \item Boundary detection: which margin inequality is active, when does it
    flip?
  \item One-sided likelihood evaluation: $\ell_j(\theta_b)$ and
    $\ell_k(\theta_b)$ to compute $\Delta U$.
  \item Event normal: total derivative $\nabla_\theta g(\theta)$ through the
    solver + filter.
  \item Reflection/refraction kernel with energy test.
\end{enumerate}

\textbf{Experiments:}
\begin{enumerate}
  \item Implement event-aware HMC with explicit reflection at value jumps.
  \item Validate: record Hamiltonian before/after crossing, check conservation
    $H^+ = H^-$ for successful refraction.
  \item Measure: crossing rate, reflection rate, acceptance rate, ESS.
  \item Compare: frozen-regime HMC (sample $\theta$ conditional on one regime
    path, then resample regime path) vs event-aware joint HMC.
\end{enumerate}

\textbf{Deliverable:} Feasibility assessment for event-aware HMC on implicit
finite-horizon OBC models. If bisection + re-solve overhead is prohibitive,
document why and pivot to alternatives (DHMC coordinate-wise, surrogate
proposal, or regime-conditional sampling).

\section{Open Questions and Future Directions}

\subsection{Is the OBC Likelihood Lipschitz Continuous?}

Even with softplus smoothing, the composed map $\theta \mapsto \log p(y \mid
\theta)$ may fail global Lipschitz continuity. The filter can amplify small
solver changes through covariance updates and log determinants. Without a
Lipschitz bound, standard HMC theory does not guarantee finite integration
error.

\textbf{Research question:} Can we derive a local Lipschitz constant as a
function of $\alpha$ (softplus temperature), filter type, and horizon? Or must
we accept that OBC models are inherently non-Lipschitz and design
integrators accordingly?

\subsection{Particle Filter Gradients Revisited}

Differentiable particle filters (LEDH-PFPF-OT) maintain diversity and should
smooth the marginal likelihood gradient. But:
\begin{itemize}
  \item Current implementation has biased score estimates;
  \item Computational cost is high (requires proposal UKF + transport);
  \item Gradient variance can still be large.
\end{itemize}

\textbf{Research question:} Is unbiased particle-filter score estimation
tractable for finite-horizon DSGE models? Can we prove that particle diversity
bounds the gradient Lipschitz constant?

\subsection{Alternative Parameterizations}

Instead of sampling $\theta$ directly, consider:
\begin{enumerate}
  \item \textbf{Regime-path augmentation.} Introduce $r_{1:T}$ as explicit
    discrete variables and sample $(\theta, r)$ jointly using mixed HMC (Zhou
    2020). The regime indicators remove the implicit boundary, but discrete
    updates and consistency checks remain.
  \item \textbf{Transformed parameters.} Reparameterize $\theta$ to align with
    the boundary geometry (e.g., polar coordinates around a regime-switching
    surface). This requires knowing the boundary shape, which is implicit.
  \item \textbf{Auxiliary-variable HMC.} Introduce latent variables that
    smooth the posterior (e.g., shadow observations, regime probabilities) and
    integrate them out after sampling.
\end{enumerate}

\textbf{Research question:} Which reparameterization, if any, makes the OBC
posterior more HMC-friendly without changing the target?

\subsection{Sequential Monte Carlo as a Likelihood Engine}

Particle MCMC (PMCMC) uses a particle filter to estimate $p(y \mid \theta)$
and treats it as a noisy but unbiased likelihood for Metropolis-within-Gibbs.
This avoids parameter-space gradients entirely.

\textbf{Tradeoffs:}
\begin{itemize}
  \item \textbf{Pro:} No gradient discontinuities, naturally handles discrete
    regime switches, can use bootstrap particle filter (simple, no transport).
  \item \textbf{Con:} Random-walk Metropolis is slow in high dimensions; no
    gradient information; particle filter must be run at every proposal.
\end{itemize}

\textbf{Research question:} For high-dimensional DSGE models, is PMCMC
competitive with event-aware HMC, or does the lack of gradient information
dominate the computational cost?

\subsection{Hybrid: Regime-Conditional HMC}

\begin{enumerate}
  \item Given a regime path $r_{1:T}$, the model is piecewise linear and the
    Kalman filter is exact.
  \item Run HMC on $p(\theta \mid y, r)$—smooth, fast, standard leapfrog.
  \item Resample $r \sim p(r \mid \theta, y)$ using a discrete sampler
    (Gibbs, Metropolis, or a regime-path proposal from the OBC solver).
  \item Alternate: HMC on $\theta$ (fixed $r$) $\to$ update $r$ (fixed
    $\theta$) $\to$ repeat.
\end{enumerate}

\textbf{Tradeoffs:}
\begin{itemize}
  \item \textbf{Pro:} Each HMC step is smooth; no event detection needed.
  \item \textbf{Con:} Regime resampling can have low acceptance if $p(r \mid
    \theta, y)$ is peaked; convergence depends on regime-parameter coupling.
\end{itemize}

\textbf{Research question:} Does regime-conditional HMC mix efficiently for
long-horizon OBC models, or does the regime sampler become a bottleneck?

\section{Interim Recommendations}

Until the research program produces validated methods, the following
recommendations apply for production DSGE inference with OBC:

\subsection{For Continuous-Kink Models (Softplus ZLB, No Integer Regimes)}

\begin{enumerate}
  \item \textbf{Characterize the geometry first.} Run Phase 1 diagnostics:
    gradient Lipschitz, crossing rate, filter variance collapse, acceptance
    degradation.
  \item \textbf{Tune aggressively.} Use small step sizes, adaptive step-size
    selection, and dense mass-matrix adaptation. Accept that mixing will be
    slower than for smooth targets.
  \item \textbf{Consider surrogate proposal.} If gradient evaluation is
    expensive, a neural surrogate (Phase 3) with hard-target correction may
    improve cost-effectiveness.
  \item \textbf{Monitor crossing diagnostics.} Report: fraction of trajectories
    crossing steep regions, acceptance conditional on crossing, ESS per
    boundary type.
\end{enumerate}

\subsection{For Value-Jump Models (Integer Regime Selector)}

\begin{enumerate}
  \item \textbf{Measure the jumps.} Phase 1 value-continuity check: are
    $\Delta \ell$ values $O(1)$ or $O(\epsilon)$?
  \item \textbf{If jumps are small:} Treat as continuous-kink and proceed as
    above.
  \item \textbf{If jumps are large:} Standard leapfrog is invalid. Options:
    \begin{itemize}
      \item Event-aware HMC with reflection/refraction (Phase 5)—expensive but
        exact.
      \item DHMC with Laplace momentum (Nishimura et al.)—simpler event
        handling, requires random order and step size.
      \item Regime-conditional HMC (hybrid sampler)—avoid joint regime-parameter
        integration.
      \item PMCMC with particle filter likelihood—no gradients, robust to
        jumps.
    \end{itemize}
  \item \textbf{Validate reversibility.} For any event-aware method, record
    reverse-replay error: trajectory $\to$ flip momentum $\to$ replay $\to$
    check distance from start.
\end{enumerate}

\subsection{For All OBC Models}

\begin{enumerate}
  \item \textbf{Separate proposal from target.} Neural approximations,
    smoothed selectors, and mollified likelihoods are allowed for proposals if
    the hard target appears in acceptance. Using them for acceptance without
    correction changes the posterior.
  \item \textbf{Freeze adaptation before inference.} If surrogate or kernel is
    adapted from chain history, freeze it before collecting posterior samples.
    Diagnostic: plot ESS over iterations; a drop at freeze indicates the
    adaptation was still improving.
  \item \textbf{Report boundary telemetry.} Every OBC inference report should
    include: geometry classification (smooth/kink/jump), boundary crossing
    rate, acceptance rate conditional on crossing, and regime occupancy
    distribution.
  \item \textbf{Compare to regime-marginalized baselines.} If feasible, run a
    particle filter or ensemble Kalman filter without HMC to check: does the
    difficulty come from the OBC or from the sampler?
\end{enumerate}

\section{Conclusion}

Occasionally-binding constraints create parameter-to-likelihood maps with
discontinuous gradients or value jumps. Standard leapfrog HMC remains
mathematically valid for continuous-kink targets but performs poorly due to
integration error at boundaries. For value-jump targets, standard leapfrog is
invalid and requires event-aware extensions (reflection, refraction, or
discontinuous coordinates).

The two-country ZLB DNS model with softplus smoothing is a continuous-kink
target. The full BGS model with integer QE regime selection may have value
jumps. Neither has been fully characterized. The proposed five-phase research
program will:
\begin{enumerate}
  \item Measure the boundary geometry and quantify gradient variation;
  \item Test event-aware integration (Tran--Kleppe truncated leapfrog);
  \item Develop surrogate-proposal HMC with hard-target correction;
  \item Explore UKF variance inflation as a gradient-smoothing heuristic;
  \item Implement reflection/refraction HMC for the hard integer-regime case.
\end{enumerate}

The interim recommendation is conservative: characterize the geometry, tune
carefully, monitor crossing diagnostics, and validate any event-aware or
surrogate method with reversibility checks and posterior agreement tests.

Particle filters remain the robust fallback for models where HMC extensions
prove intractable. But their current limitations (biased gradients, high cost)
motivate continued work on making gradient-based MCMC viable for OBC models.
This chapter documents the problem and the research path. The solution is not
yet complete.
